The same determinant-open quiver governs the linear pair
and the scalar extension of its unitary counterpart. We first establish
its closed-orbit geometry and normal spaces over $F$. The
unitary application will retain its rational descent data and the
additional multiplier direction.
The algebraic statements and constructions of this section hold over
every field of characteristic zero; in Section~\ref{sec:unitary-geometry}
we also apply them over $E$ and $\Fb$.

\subsection{The determinant open and its rational closed orbits}\label{m1:sec:geometry}
\subsubsection{Determinant-relative stability}\label{m1:subsec:determinant-stability}
Let $\qvr$ be the quiver with vertices $U,Q,P$ and arrows
$A:U\to P$, $B:Q\to P$, $C:U\to Q$, and $D:Q\to Q$. Dimension vectors
are ordered as $(\dim U,\dim Q,\dim P)$. Its representation space in
dimension $(2,2,2)$ is
\[
 V=\Hom(U,P)\oplus\Hom(Q,P)\oplus\Hom(U,Q)\oplus\End(Q),
\]
with base-change group $L=\GL(P)\times\GL(Q)\times\GL(U)$.
For $x=(A,B,C,D)$ its action, the transpose involution, and the
induced automorphism of $L$ are
\begin{align}
 (p,q,u)x&=(pAu^{-1},pBq^{-1},qCu^{-1},qDq^{-1}),
 \label{m1:eq:residual-action}\\*
 \tau x&=(A^t,C^t,B^t,D^t),\qquad
 \alpha(p,q,u)=(u^{-t},q^{-t},p^{-t}).\label{m1:eq:tau}
\end{align}
 For $E=(U,Q,P;A,B,C,D)$ put
\begin{equation}\label{m1:eq:block-map}
 \Phi_E:U\oplus Q\longrightarrow P\oplus Q,
 \qquad\Phi_E(u,q)=(Au+Bq,Cu+Dq).
\end{equation}
The residual variety is the principal open
\begin{equation}\label{m1:eq:X}
 \cX=\{E\in V:f(E)\neq0\},\qquad f(E)=\det\Phi_E.
\end{equation}
The determinant is a semi-invariant:
\begin{equation}\label{m1:eq:det-character}
 f((p,q,u)E)=\chi_{\det}(p,q,u)f(E),
 \qquad\chi_{\det}(p,q,u)=\det p/\det u.
\end{equation}
For a field $k$ and a character $\gamma$ with values in $k^\times$,
write $k_\gamma$ for the one-dimensional representation $g\cdot z=\gamma(g)z$,
and write $k_{\triv}$ for the trivial representation.

A subrepresentation $N=(U',Q',P')$ satisfies
$AU'\subset P'$, $BQ'\subset P'$, $CU'\subset Q'$, and $DQ'\subset Q'$.
Define its balance by
\begin{equation}\label{m1:eq:theta}
 \theta(N)=\dim P'-\dim U'.
\end{equation}
We call $N$ balanced when $\theta(N)=0$.

\begin{lemma}\label{m1:lem:theta-nonnegative}
If $E\in\cX$, then $\theta(E)=0$ and $\theta(N)\geq0$ for every
subrepresentation $N$. Every balanced subrepresentation and its quotient
have invertible induced block maps.
\end{lemma}
\begin{proof}
The restriction of $\Phi_E$ is an injection
$U'\oplus Q'\to P'\oplus Q'$, so $\dim U'\leq\dim P'$. If equality
holds, this restriction is an isomorphism, and so is the induced quotient
map. The argument applies successively to balanced filtration factors.
\end{proof}

Let $\cAZero$ be the full subcategory of finite-dimensional quiver
representations with $\theta(E)=0$ and $\theta(N)\geq0$ for all
subrepresentations. A nonzero object is stable if the latter inequality
is strict for every nonzero proper subrepresentation, and polystable if
it is a direct sum of stable objects. This is the zero-slope semistable
category, with slope $\theta(E)/(\dim U+\dim Q+\dim P)$ for nonzero $E$.

\begin{lemma}\label{m1:lem:A0-abelian}
The category $\cAZero$ is abelian of finite length. Its simple objects
are precisely its stable objects.
\end{lemma}
\begin{proof}
For a morphism $h:E\to E'$ in $\cAZero$, its kernel and image have
nonnegative balance. Additivity and $\theta(E)=0$ force both balances
to be zero; the cokernel is then balanced as well. Nonnegativity passes
to subobjects of the kernel and image, and to the cokernel by lifting
subobjects to $E'$. These are therefore the kernel, image, and cokernel
in the full subcategory. Direct sums satisfy the same condition by
intersecting a subobject with one summand and projecting to the other.
Every balanced subrepresentation of an object of $\cAZero$ belongs to
$\cAZero$, so simplicity is exactly stability. Total vector-space
dimension strictly decreases in a proper subobject or quotient, proving
finite length.
\end{proof}

Over an algebraically closed field, these definitions agree with
character-GIT semistability and stability for $\chi_{\det}$ by
\cite[Proposition~3.1]{King94}. In the semistable locus, closed orbits
correspond to polystable objects by \cite[Proposition~3.2]{King94}.
The semistable locus also contains points with singular block map,
such as an $A$-only stable factor together with a $Q$-supported
zero-loop factor. We prove the closed-orbit criterion on $\cX$ directly.

\begin{proposition}[Polystability on the determinant open]
\label{m1:prop:polystability}
For $E\in\cX$, its geometric $L$-orbit is closed in $\cX$ if and only if
$E_{\Fb}$ is semisimple in $\cAZero$, equivalently polystable.
\end{proposition}
\begin{proof}
Work over $\Fb$. The map $v\mapsto(v,f(v)^{-1})$ identifies $\cX$
equivariantly with the closed invariant graph
\begin{equation}\label{m1:eq:closed-graph}
 \widehat\cX=\{(v,z)\in V\oplus\Fb_{\chi_{\det}^{-1}}:zf(v)=1\}.
\end{equation}
We apply \cite[Proposition~2.6(i)]{King94}, with the trivial character,
to this ambient linear representation of the reductive group $L$.
For that character the semistable locus is the entire representation
space, and its pairing with every one-parameter subgroup is zero.
Thus an orbit is closed if and only if every existing one-parameter
limit stays in that orbit. Since the graph is closed, any such limit
from the graph remains in the graph.

If $E$ is not semisimple, choose a Jordan--H\"older filtration
$0=E_0\subset E_1\subset\cdots\subset E_r=E$ in $\cAZero$, with
$S_i=E_i/E_{i-1}$. Split the vertex filtrations as vector spaces, and
give the $i$th graded part at every vertex a weight $w_i$, where
$w_1>\cdots>w_r$. For any arrow, its block from the $j$th source part
to the $i$th target part is zero unless $i\leq j$, and is multiplied by
$t^{w_i-w_j}$. The one-parameter limit is therefore $\bigoplus_i S_i$.
Its determinant weight is
$\sum_i w_i\theta(S_i)=0$, so the graph coordinate stays fixed and the
limit lies in $\cX$. It is not isomorphic to $E$, proving nonclosedness.

Conversely suppose $E$ is semisimple in $\cAZero$, and a one-parameter
subgroup has a limit in \eqref{m1:eq:closed-graph}. Its weights decompose
every vertex space; let $F^cE$ consist of the weights at least $c$.
An arrow block from weight $b$ to weight $a$ is multiplied by $t^{a-b}$;
existence of the limit forces it to vanish for $a<b$. Thus every $F^cE$
is a subrepresentation. If the distinct weights are
$m_1<\cdots<m_s$, put $\delta_i=\dim P_{m_i}-\dim U_{m_i}$. Then
$\sum_i\delta_i=0$ and $\theta(F^{m_i}E)=\sum_{j\geq i}\delta_j\geq0$.
The graph equation forces the determinant weight to be zero: both
$f(\lambda(t)E)=t^k f(E)$ and its inverse have finite nonzero limits,
so $k=0$. Summation by parts gives
\begin{equation}\label{m1:eq:balanced-weight-identity}
 0=\sum_i m_i\delta_i
   =\sum_{i=2}^s(m_i-m_{i-1})\theta(F^{m_i}E).
\end{equation}
Each term is nonnegative, hence every threshold subobject is balanced
and belongs to $\cAZero$. This filtration splits because $E$ is
semisimple there. Its associated graded, which is the one-parameter
limit, is therefore isomorphic to $E$. The affine criterion proves
closedness.
\end{proof}
For $E$ defined over $F$, geometric closedness is
equivalent to analytic closedness of $L(F)E$ by
\cite[Theorem~2.3.8]{AGDescent}; all rational orbit identifications below
are proved over $F$.

\subsubsection{Stable factors and arrow ranks}\label{m1:subsec:stable-factors}
Duality interchanging $U$ and $P$ preserves stability. Indeed a proper
nonzero subobject of the dual is dual to a quotient $E/N$, and its
balance is $-\theta(E/N)=\theta(N)>0$ for stable $E$.

\begin{lemma}\label{m1:lem:stable-factors}
Over $\Fb$, a stable balanced factor of a point of $\cX$, of dimension
$(r,s,r)$, has $s=1$ if $r=0$, and $r=1$ if $s=0$. If $r,s>0$, its
$C$ is injective and its $B$ is surjective. The only proper stable
dimension vectors in ambient dimension $(2,2,2)$ are
\[
 (1,0,1),\quad(0,1,0),\quad(1,1,1),\quad(1,2,1).
\]
\end{lemma}
\begin{proof}
For $r=0$ only the loop remains, and a loop on a space of dimension
greater than one over $\Fb$ has a proper invariant subspace. For $s=0$,
the block map is the isomorphism $A:U\to P$; a proper subspace $U'$
would give a balanced subobject $(U',0,AU')$. These prove the first two
assertions. If $r,s>0$ and $U_0=\ker C\neq0$, injectivity of $\Phi_E$
makes $A|_{U_0}$ injective. Then $(U_0,0,AU_0)$ is a proper nonzero
balanced subobject, contradicting stability. Duality gives surjectivity
of $B$. Hence $r\leq s$, which gives the list under the ambient
dimension bounds. The same argument applies to a full-dimensional
stable object.
\end{proof}

Denote the first two factors by $\Efac$ and $\Qfac_e$. On $\Efac$ we
normalize $A=1$, and on $\Qfac_e$ the loop is $e\neq0$. A stable
$(1,1,1)$-factor is normalized as
\begin{equation}\label{m1:eq:R-factor}
 \Rfac_{x,d}:\qquad \Phi=\begin{pmatrix}x&1\\1&d\end{pmatrix},
 \qquad xd\neq1.
\end{equation}
It is stable because $B$ and $C$ are nonzero, leaving no proper balanced
subobject. We denote a stable $(1,2,1)$-factor by $\Hfac$.

\begin{corollary}\label{m1:prop:rank-exclusion}
Every closed orbit in $\cX$ has $\rk B=\rk C$. Its rank type is
$(2,2)$, $(1,1)$, or $(0,0)$.
\end{corollary}
\begin{proof}
On each stable factor the two ranks agree: they are zero on $\Efac$
and $\Qfac_e$, one on $\Rfac$ and $\Hfac$, and two on a full stable
factor. Ranks add under the polystable decomposition of
Proposition~\ref{m1:prop:polystability}. This excludes all six unequal-rank
types, including those with only one nonzero arrow.
\end{proof}
Geometrically the rank-zero factors are $\Efac^{\oplus2}$ and a
semisimple two-dimensional loop object. Rank one has either
$\Efac\oplus\Rfac\oplus\Qfac$ or $\Efac\oplus\Hfac$. Rank two is
stable of full dimension or a sum of two $(1,1,1)$-factors, possibly
isomorphic. This is a geometric list; a nonsplit rational simple may
split into geometric factors and need not have one of the proper
dimension vectors above.

\subsubsection{Dense pairs and rational semisimplicity}\label{m1:subsec:dense-geometry}
The regular $L$-invariant function
$\rho(E)=\det B\det C/\det\Phi_E$ shows that an orbit in the dense
rank-$(2,2)$ locus cannot degenerate to a lower rank in $\cX$: its
closure lies in the nonzero fiber $\rho^{-1}(\rho(E))$.
Taking $p=qB^{-1}$ and $u=qC$ gives the gauge $B=C=I$ with
\begin{equation}\label{m1:eq:dense-gauge}
 X=B^{-1}AC^{-1},\qquad Y=D,
 \qquad (X,Y)\mapsto(qXq^{-1},qYq^{-1}).
\end{equation}
The condition is $\Delta(X,Y)=\det\left(\begin{smallmatrix}X&I\\I&Y\end{smallmatrix}\right)\neq0$.
As $\Delta$ is conjugation-invariant, a pair orbit's closure in $M_2^2$
also remains in this open.

\begin{proposition}\label{m1:prop:dense-closed}
A dense rational orbit is closed exactly when
$M=(F^2;X,Y)$ is semisimple as an $F\langle x,y\rangle$-module.
Its four rational types and stabilizers are: an absolutely irreducible
module, with stabilizer $F^\times$; two distinct rational characters,
with stabilizer $(F^\times)^2$; a rational irreducible module splitting
over a quadratic field $K/F$, with stabilizer $K^\times$; and two equal
characters, with stabilizer $\GL_2(F)$.
\end{proposition}
\begin{proof}
In the gauge, a balanced subobject has $\dim U'=\dim P'$ and
$U'\subset Q'\subset P'$, because $B=C=I$. Thus all three spaces agree.
Such subobjects are exactly the common $X,Y$-invariant subspaces.
This identifies geometric polystability with semisimplicity of the
free-algebra module by Proposition~\ref{m1:prop:polystability}.

The rational comparison requires both directions. If $M$ is semisimple
over $F$, let $A_M\subset\End_F(M)$ be the finite-dimensional image of
$F\langle x,y\rangle$. Its radical annihilates the faithful semisimple
module $M$, hence is zero. In characteristic zero, a finite-dimensional
semisimple algebra is separable: its centers are finite separable
extensions and its central simple factors split over $\Fb$. Therefore
$A_M\otimes_F\Fb$ and $M\otimes_F\Fb$ are semisimple. Conversely, for
an $F$-defined submodule $W\subset M$, a retraction over $\Fb$ satisfies
linear equations over $F$: it commutes with $X,Y$ and is the identity
on $W$. Row reduction gives a solution over $F$. Thus geometric
semisimplicity implies rational semisimplicity.

The stabilizer is $\End_{A_M}(M)^\times$. Absolute irreducibility gives
$\End_{A_M}(M)=F$. Two distinct rational characters give $F\oplus F$,
and two equal characters give $M_2(F)$. In the remaining irreducible
case, $D_0=\End_{A_M}(M)$ is a division algebra, and $\dim_FD_0$ divides
$\dim_FM=2$. If $D_0=F$, the finite-dimensional density theorem gives
$A_M=M_2(F)$, contrary to failure of absolute irreducibility. Hence
$\dim_FD_0=2$; any nonscalar element generates it, so $D_0$ is a
quadratic field $K$. Identifying $M$ with $K$, the pair consists of
multiplication operators $m_x,m_y$, at least one nonscalar. These exhaust
all semisimple modules of dimension two.
\end{proof}

\subsubsection{Rank zero and rank one over \texorpdfstring{$F$}{F}}
\label{m1:subsec:rational-boundary}
If $B=C=0$, invertibility gives $A,D\in\GL_2(F)$ and base change
normalizes $A=I$. The object is $\Efac^{\oplus2}\oplus\Qfac_D$, where
$\Qfac_D$ is supported on the loop vertex. Polystability says exactly
that $D$ is semisimple. Rational canonical form yields the following.

\begin{proposition}\label{m1:prop:rankzero}\label{m1:prop:rank00-classification}
Every closed rank-zero orbit has a representative
$x_D=(I,0,0,D)$ with $D\in\GL_2(F)$ semisimple. Its rational type is
scalar $D=dI$ with $d\neq0$, split regular with two distinct nonzero
eigenvalues, or nonsplit regular with irreducible quadratic characteristic
polynomial.
\end{proposition}

For rank one, $U_0=\ker C$ and $P_0=AU_0$ are rational lines, and
$\Efac=(U_0,0,P_0)$ is a balanced subrepresentation.

\begin{lemma}\label{m1:lem:rank11-split}
If $E$ has closed orbit and rank type $(1,1)$, there is an $F$-rational
splitting $E\simeq\Efac\oplus\mathcal H$, with
$\dim\mathcal H=(1,2,1)$.
\end{lemma}
\begin{proof}
The inclusion of $\Efac$ splits over $\Fb$ by polystability. A retraction
is a triple of linear maps satisfying the quiver morphism equations and
restricting to the identity on $\Efac$. All these equations are linear
over $F$, so solvability over $\Fb$ implies solvability over $F$. Its
kernel is the required complement, geometrically semisimple in
$\cAZero$.
\end{proof}
Write the complement's arrows as $c:F\to Q$, $b:Q\to F$, $D:Q\to Q$,
and its remaining scalar arrow as $a$. Both $b$ and $c$ are nonzero.
Put $L_c=\operatorname{im}c$ and $K_b=\ker b$.

\begin{lemma}\label{m1:lem:rank11-balanced}
The only proper nonzero balanced subobjects of $\mathcal H$ are
$(0,K_b,0)$, present exactly when $K_b$ is $D$-stable, and
$(F,L_c,F)$, present exactly when $L_c$ is $D$-stable. Thus $\mathcal H$
is geometrically stable exactly when neither line is $D$-stable.
\end{lemma}
\begin{proof}
The two end dimensions of a balanced subobject are either both zero or
both one. In the first case its middle space is a nonzero subspace of
$K_b$, hence is $K_b$. In the second it must contain $L_c$, and
properness forces it to be $L_c$. Each is a subobject precisely when
its middle line is $D$-stable. This argument holds over every scalar
extension, and both candidate lines are defined over $F$.
\end{proof}

\begin{proposition}[Rank-one dichotomy]\label{m1:prop:rank11}
A rank-one orbit is closed if and only if it admits the rational
splitting above and either neither of $L_c,K_b$ is $D$-stable, or both
are $D$-stable and distinct. The first case is
$E=\Efac\oplus\Hfac$, with $\Hfac$ stable. The second is
$E=\Efac\oplus\Rfac\oplus\Qfac$.
\end{proposition}
\begin{proof}
The splitting is necessary by Lemma~\ref{m1:lem:rank11-split}. If neither
line is stable, Lemma~\ref{m1:lem:rank11-balanced} proves stability of the
complement over $\Fb$. If both are stable and distinct, then
$Q=L_c\oplus K_b$ is a rational $D$-stable decomposition, with $b$
nonzero on $L_c$. Its two summands have stable dimension vectors
$(1,1,1)$ and $(0,1,0)$, proving polystability.
If exactly one line is stable, or both are stable and coincide, the
balanced subobject supplied by the lemma has no complement: a complement
to $(0,K_b,0)$ must use $L_c$, and a complement to $(F,L_c,F)$ must use
$K_b$. Either splitting forces the lines to be stable and distinct.
Thus these remaining cases are not semisimple, and their orbits are
not closed by Proposition~\ref{m1:prop:polystability}.
\end{proof}

The decomposable family has the representatives
\begin{equation}\label{m1:eq:decomposable-rep}
 A=\diag(x,1),\quad B=C=E_{11},\quad D=\diag(d,e),
 \qquad e\neq0,\quad xd\neq1.
\end{equation}
These are exactly the invertibility conditions on
$\Efac\oplus\Rfac_{x,d}\oplus\Qfac_e$; they allow $x=0$ and $d=e$.

For the stable-system family choose the rational cyclic basis $(c,Dc)$
of $Q$, possible because $L_c$ is not $D$-stable. Then $c=e_1$ and
$D=\left(\begin{smallmatrix}0&-d\\1&s\end{smallmatrix}\right)$.
Adjoining $\Efac$ gives
\begin{equation}\label{m1:eq:stable-rep-general}
 A=\diag(a,1),\quad
 B=\begin{pmatrix}b_0&b_1\\0&0\end{pmatrix},\quad C=E_{11},\quad
 D=\begin{pmatrix}0&-d\\1&s\end{pmatrix}.
\end{equation}
The determinant on the $(1,2,1)$-factor is
$\Delta=ad-sb_0+b_1\neq0$. Scaling the first $P$-coordinate by
$\Delta^{-1}$ scales $a,b_0,b_1,\Delta$ together and fixes $C,D$ and
the complementary identity block. Relabeling the scaled parameters gives
\begin{equation}\label{m1:eq:b1-normalized}
 b_1=1-ad+sb_0.
\end{equation}
Writing $b=(b_0,b_1)$, set
\begin{equation}\label{m1:eq:H}
 H=\begin{pmatrix}b_0&b_1\\b_1&sb_1-db_0\end{pmatrix}
   =\begin{pmatrix}b\\bD\end{pmatrix},\qquad\det H\neq0.
\end{equation}
The rows $b,bD$ are dependent exactly when $\ker b$ is $D$-stable:
dependence means $bD$ vanishes on $\ker b$, and conversely any functional
vanishing there is a multiple of the nonzero $b$. Thus \eqref{m1:eq:H} is
precisely the remaining stability condition. These normalizations also apply to nilpotent and repeated-root values of $D$.

\begin{theorem}[Rational closed-orbit classification]\label{m1:thm:closed-orbits}
Every closed $L(F)$-orbit in $\cX(F)$ is represented by exactly one of
the following families of types: a dense semisimple pair from
Proposition~\ref{m1:prop:dense-closed}; a rank-zero point from
Proposition~\ref{m1:prop:rankzero}; a decomposable representative
\eqref{m1:eq:decomposable-rep}; or a stable-system representative
\eqref{m1:eq:stable-rep-general} satisfying
\eqref{m1:eq:b1-normalized}--\eqref{m1:eq:H}. Each representative above has a closed orbit. The nine families are disjoint; within each family,
the parameters may represent the same orbit.
\end{theorem}
\begin{proof}
Corollary~\ref{m1:prop:rank-exclusion} gives the three ranks, and the dense,
rank-zero, and rational rank-one classifications exhaust them. Their
proofs establish both closedness and the exhaustion of the rational orbits.
\end{proof}

\subsection{Normal representations and rational transpose}\label{m1:sec:normals}
The classification supplies the objects at which descent is applied.
We now identify their stabilizers, normal representations, and
corrected transpose operators over $F$.

\subsubsection{The Hom--Ext complex}\label{m1:subsec:hom-ext}
For representations $S,T$ of $\qvr$ over $F$, let
\[
 \begin{gathered}
 C^0(S,T)=\bigoplus_v\Hom_F(S_v,T_v),\qquad
 C^1(S,T)=\bigoplus_{a:v\to w}\Hom_F(S_v,T_w),\\
 d(h)_a=h_wS_a-T_ah_v.
 \end{gathered}
\]

\begin{lemma}\label{m1:lem:hom-ext}
There is a natural exact sequence
\begin{equation}\label{m1:eq:hom-ext-complex}
 0\longrightarrow\Hom_{\qvr}(S,T)\longrightarrow C^0(S,T)
 \xrightarrow{d}C^1(S,T)\longrightarrow\Ext^1_{\qvr}(S,T)
 \longrightarrow0.
\end{equation}
For $E\in\cX$ this identifies, equivariantly under $L_E=\Aut_{\qvr}(E)$,
\begin{equation}\label{m1:eq:full-normal-ext}
 N_E=T_E\cX/T_E(LE)\simeq\Ext^1_{\qvr}(E,E),
 \qquad\dim_F N_E=4+\dim_F\End_{\qvr}(E).
\end{equation}
\end{lemma}
\begin{proof}
The kernel consists exactly of quiver morphisms. Split an extension
$0\to T\to Z\to S\to0$ at each vertex as vector spaces. Its arrow
matrices are $\left(\begin{smallmatrix}T_a&c_a\\0&S_a\end{smallmatrix}\right)$.
There are no relations on the arrows, including the loop, so $c_a$ is
arbitrary. Changing the vertex splittings changes $(c_a)$ by an element
of $\operatorname{im}d$, up to sign. Conversely this change gives an
equivalent extension. Addition and scalar multiplication of the upper
right blocks give the corresponding operations on extension classes,
proving the cokernel assertion.

Since $\cX$ is open, $T_E\cX=C^1(E,E)$. Its orbit differential is
\begin{equation}\label{m1:eq:orbit-differential}
 d_E(P_0,Q_0,U_0)
 =(P_0A-AU_0,\ P_0B-BQ_0,\ Q_0C-CU_0,\ Q_0D-DQ_0).
\end{equation}
Its kernel is $\End_{\qvr}(E)$ and its image is $T_E(LE)$.
Taking the quotient proves \eqref{m1:eq:full-normal-ext}, including the
dimension formula $16-12+\dim\End_{\qvr}(E)$. All constructions commute
with automorphisms of $E$.
\end{proof}
For dimension vectors $\mathbf d=(u,q,p)$ and
$\mathbf e=(u',q',p')$, the Euler form from
\eqref{m1:eq:hom-ext-complex} is
\[
 \langle\mathbf d,\mathbf e\rangle
 =uu'+qq'+pp'-up'-qp'-uq'-qq'.
\]
On balanced vectors it is symmetric:
\begin{equation}\label{m1:eq:balanced-euler}
 \langle(r,s,r),(r',s',r')\rangle=rr'-sr'-rs',\qquad
 \dim\Ext^1(S,T)=\dim\Hom(S,T)-\langle\dim S,\dim T\rangle.
\end{equation}
Distinct stable objects have no Hom between them, since they are
nonisomorphic simples in the full subcategory $\cAZero$. For an
absolutely stable object defined over $F$, its endomorphism algebra is
$F$, by Schur's lemma after scalar extension and descent of the linear
morphism equations. Thus for pairwise nonisomorphic absolutely stable
$S_i$ defined over $F$ and finite-dimensional $F$-spaces $M_i$,
\begin{equation}\label{m1:eq:multiplicity-spaces}
 \begin{split}
 E=\bigoplus_i S_i\otimes_F M_i\quad&\Longrightarrow\quad
 \Aut(E)=\prod_i\GL(M_i),\\
 \Ext^1(E,E)&=\bigoplus_{i,j}
 \Hom_F(M_i,M_j)\otimes_F\Ext^1(S_i,S_j).
 \end{split}
\end{equation}
These identities follow by decomposing the complex into its $(i,j)$
blocks. On the factor $\Hom_F(M_i,M_j)$ in \eqref{m1:eq:multiplicity-spaces}, the action is $f\mapsto g_jfg_i^{-1}$.
The nonsplit rational cases below are computed in their quadratic-field models.

\subsubsection{Normal representations for dense orbits}\label{m1:subsec:dense-normals}
At the gauge $B=C=I$ the differential of \eqref{m1:eq:dense-gauge} is
\[
 (\dot A,\dot B,\dot C,\dot D)\longmapsto
 (\dot A-\dot B X-X\dot C,\dot D).
\]
It is onto $M_2^2$. Its kernel consists of orbit directions: use
$Q_0=0$, $P_0=\dot B$, $U_0=-\dot C$ in
\eqref{m1:eq:orbit-differential}. The image of a general orbit direction
is $([Q_0,X],[Q_0,Y])$. Hence
\begin{equation}\label{m1:eq:dense-normal}
 N_{X,Y}\simeq\Ext^1_{F\langle x,y\rangle}(M,M)
 =\frac{M_2(F)\oplus M_2(F)}
 {\{([Z,X],[Z,Y]):Z\in M_2(F)\}}.
\end{equation}
The Ext group is taken over the free algebra, so both generators may
vary. The two-loop Euler form is
$\langle m,n\rangle=-mn$.

\begin{lemma}[Dense rational transporters]\label{m1:lem:dense-transpose}
Every closed dense pair has an $F$-rational symmetric invertible matrix
$H_0$ with $H_0X=X^tH_0$ and $H_0Y=Y^tH_0$. The transporter in $L$ is
$(H_0^{-1},H_0^{-1},H_0^{-1})$, and the corrected operator on
\eqref{m1:eq:dense-normal} is induced by
\begin{equation}\label{m1:eq:dense-transpose}
 (P_1,Q_1)\longmapsto(H_0^{-1}P_1^tH_0,H_0^{-1}Q_1^tH_0).
\end{equation}
\end{lemma}
\begin{proof}
For an absolutely irreducible pair, the space of symmetric matrices has
dimension three over $F$. Each of $H_0X-X^tH_0$ and $H_0Y-Y^tH_0$ is
skew-symmetric, imposing at most one linear condition. There is a
nonzero rational symmetric solution. If it were singular, its kernel
would be a common invariant line, contradicting irreducibility. For
split and scalar pairs choose diagonal representatives and $H_0=I$.
For a nonsplit pair on $K$, take the Gram matrix of
$\operatorname{Tr}_{K/F}(uv)$; multiplication operators are self-adjoint
for this nondegenerate symmetric form. To return from a normal-form pair
$(X_0,Y_0)$ to $(X,Y)=g(X_0,Y_0)g^{-1}$, replace its intertwiner
$H_{\mathrm{rep}}$ by $H_0=g^{-t}H_{\mathrm{rep}}g^{-1}$.
This matrix is symmetric and invertible and satisfies
$H_0X=X^tH_0$ and $H_0Y=Y^tH_0$.
The transporter $(H_0^{-1},H_0^{-1},H_0^{-1})$ carries the transposed point back to the original one. Differentiation then gives \eqref{m1:eq:dense-transpose}.
Symmetry makes the corrected operator square to one.
\end{proof}

\begin{proposition}[Normal space for an absolutely irreducible pair]\label{m1:prop:abs-irred-normal}
For an absolutely irreducible pair, the differentials of
\[
 f_1=\tr X,\quad f_2=\det X,\quad f_3=\tr Y,\quad
 f_4=\det Y,\quad f_5=\tr(XY)
\]
identify the normal space with $F^5$. Both the scalar stabilizer and the
corrected transpose act trivially on it.
\end{proposition}
\begin{proof}
For a variation $(P_1,Q_1)$, the trace and product differentials are
$\tr P_1$, $\tr Q_1$, and $\tr(P_1Y+XQ_1)$. Since
$\operatorname{adj}(X)=(\tr X)I-X$, the first two differentials span
$\tr P_1$ and $\tr(XP_1)$, and likewise the third and fourth span
$\tr Q_1$ and $\tr(YQ_1)$. A linear relation among all five would give
\[
 a\tr P_1+b\tr(XP_1)+c\tr Q_1+d\tr(YQ_1)
 +e\tr(P_1Y+XQ_1)=0.
\]
Nondegeneracy of the trace pairing, after separating the variables,
implies $aI+bX+eY=0$ and $cI+dY+eX=0$. If $e\neq0$, then
$Y\in F[I,X]$, so an eigenline of $X$ over $\Fb$ is also $Y$-stable,
contradicting absolute irreducibility. Thus $e=0$. Neither matrix can
be scalar, since the other has an eigenline over $\Fb$; hence
$a=b=c=d=0$. The differentials are independent.

They vanish on the orbit tangent, whose dimension is $4-1=3$. Their
common kernel has dimension $8-5=3$, and is therefore exactly that
tangent. This agrees with the self-Ext dimension $1+2^2=5$ and gives
explicit coordinates on the normal space. The map
\eqref{m1:eq:dense-transpose} fixes the base point and all five functions.
Differentiation shows that it fixes every normal coordinate. The
stabilizer is central and acts trivially.
\end{proof}

For the split type write
$X=\diag(x_1,x_2)$, $Y=\diag(y_1,y_2)$, with
$(x_1,y_1)\neq(x_2,y_2)$. The two character factors have self-Ext
dimension two and cross-Ext dimension one in each direction. Thus, with
$\alpha(t_1,t_2)=t_1/t_2$,
\begin{equation}\label{m1:eq:dense-split-normal}
 S=(F^\times)^2,\qquad N_{X,Y}=F^4_{\triv}\oplus F_\alpha
 \oplus F_{\alpha^{-1}}.
\end{equation}
The four diagonal parameter variations give the self-Ext coordinates and
are fixed by transpose. On the off-diagonal quotient transpose
interchanges the two one-dimensional spaces. Its square is one, so a
basis in one and its image in the other make this a literal swap over
$F$. The character is rationally surjective, since $\alpha(t,1)=t$.

For the nonsplit type, identify $M=K$, write $m_z$ for multiplication
by $z\in K$, and write $z\mapsto\bar z$ for the nontrivial automorphism
of $K/F$. Then $X=m_x,Y=m_y$, and with $\iota(w)=\bar w$ there is a rational decomposition
$\End_F(K)=K\oplus K\iota$. The two $K$-summands in the pair of
endomorphism spaces give four trivial directions. On the remaining
part,
\[
 [m_d\iota,m_x]=m_{d(\bar x-x)}\iota,\qquad
 [m_d\iota,m_y]=m_{d(\bar y-y)}\iota.
\]
Put $a=\bar x-x$, $b=\bar y-y$. At least one is nonzero, and hence
\begin{equation}\label{m1:eq:dense-nonsplit-quotient}
 N_{X,Y}=F^4_{\triv}\oplus V_{\mathrm{ns}},\qquad
 V_{\mathrm{ns}}=K^2/K(a,b),\qquad
 \ell(d_1,d_2)=bd_1-ad_2.
\end{equation}
The kernel of $\ell$ is $K(a,b)$, giving $V_{\mathrm{ns}}\simeq K$.
Conjugation by $t\in K^\times$ multiplies both root coordinates by
$t/\bar t$. For the trace-form adjoint,
\[
 H_0^{-1}(m_c+m_d\iota)^tH_0=m_c+m_{\bar d}\iota,
\]
because $\operatorname{Tr}(d\bar u v)=\operatorname{Tr}(\bar d u\bar v)$.
Thus the adjoint fixes the four multiplication directions, conjugates
the two ambient root coordinates, and preserves $K(a,b)$. Since
$\bar a=-a$ and $\bar b=-b$, its action after $\ell$ is
\begin{equation}\label{m1:eq:dense-nonsplit-action}
 t:d\longmapsto(t/\bar t)d,\qquad \tau_E:d\longmapsto-\bar d.
\end{equation}
The effective group is $K^\times/F^\times\simeq K^1$, by Hilbert 90.

For the scalar pair $X=xI,Y=yI$, let $\chi$ be the one-dimensional
module on which the two generators act by $x$ and $y$. The orbit tangent in
\eqref{m1:eq:dense-normal} is zero. Thus $S=\GL_2(F)$ acts by simultaneous
conjugation on $N_{X,Y}=M_2(F)^2$, and transpose is componentwise. In
multiplicity-space terms this is $M_2(F)\otimes_F\Ext^1(\chi,\chi)$,
where $\dim\Ext^1(\chi,\chi)=2$.

\subsubsection{Normal representations for rank-zero orbits}\label{m1:subsec:rankzero-normal}
At $x_D=(I,0,0,D)$, the stabilizer is
\[
 S_D=\{(g,t,g):g\in\GL_2(F),\ t\in Z_{\GL_2(F)}(D)\}
 \simeq\GL_2(F)\times Z_{\GL_2(F)}(D).
\]
This also follows from $\Efac^{\oplus2}\oplus\Qfac_D$: there are no
Hom terms between the two supports. The loop centralizer is
$(F^\times)^2$, $K^\times$, or $\GL_2(F)$ for split regular, nonsplit
regular, or scalar $D$, respectively. Put
$\mathfrak z(D)=\{Z\in M_2(F):ZD=DZ\}$. The differential specializes to
$(P_0-U_0,0,0,[Q_0,D])$, so
\begin{equation}\label{m1:eq:rankzero-normal}
 N_D=M_2(F)\oplus M_2(F)\oplus M_2(F)/[M_2(F),D]
 \simeq M_2(F)^2\oplus\mathfrak z(D).
\end{equation}
Semisimplicity of $\operatorname{ad}D$ gives the direct sum
$M_2(F)=[M_2(F),D]\oplus\mathfrak z(D)$ over $F$. Equivalently, the
trace pairing is nondegenerate on the centralizer. The action is
\begin{equation}\label{m1:eq:rankzero-action}
 (g,t)(X,Y,Z)=(gXt^{-1},tYg^{-1},tZt^{-1}).
\end{equation}

For regular $D$, choose a symmetric matrix $S$ such that $SD=D^tS$
and $St=t^tS$ for every $t\in F[D]^\times$. In the split eigenbasis use
$S=I$; in the nonsplit model $K=F[D]$, use the trace form. The rational
transporter $(I,S^{-1},I)$ corrects the point and induces
\begin{equation}\label{m1:eq:rankzero-transpose}
 \tau_D(X,Y,Z)=(Y^tS,S^{-1}X^t,S^{-1}Z^tS).
\end{equation}
This squares to one and conjugates $(g,t)$ to $(g^{-t},t^{-1})$.
Every $Z=aI+bD$ is fixed by its last component, since
$S^{-1}Z^tS=aI+bD$. The torus acts trivially there as well. Thus the
centralizer directions are $F^2$ fixed by both actions. As a concrete
alternative, for the companion matrix
$D=\left(\begin{smallmatrix}0&-d\\1&s\end{smallmatrix}\right)$,
$j=\left(\begin{smallmatrix}-s&1\\1&0\end{smallmatrix}\right)$ satisfies
$jD^t=Dj$; the trace-form construction above also determines the action
on both legs.

For scalar $D=dI$, with $d\neq0$, no correction is needed:
\begin{equation}\label{m1:eq:rankzero-scalar}
 S_D=\GL_2(F)^2,\quad N_D=M_2(F)^2\oplus\mathfrak{gl}_2(F),\quad
 \tau_D(X,Y,Z)=(Y^t,X^t,Z^t).
\end{equation}
The last summand carries the adjoint action. The
regular cases will use the Two-Leg Lemma with two fixed parameters;
the scalar case uses the Identity Descendant on all three summands.

\subsubsection{The decomposable case}\label{m1:subsec:decomposable-normal}
At \eqref{m1:eq:decomposable-rep}, the three factors
$\Efac,\Rfac_{x,d},\Qfac_e$ are pairwise nonisomorphic and absolutely
stable. Their scalar automorphisms give the embedding
\begin{equation}\label{m1:eq:decomposable-stabilizer}
 S=\{(\diag(r,b),\diag(r,s),\diag(r,b)):r,b,s\in F^\times\}.
\end{equation}
Their dimension vectors differ, so neither $x=0$ nor $d=e$ merges
factors or enlarges this stabilizer. Put
\begin{equation}\label{m1:eq:decomposable-characters}
 \alpha(r,b,s)=r/s,\qquad\beta(r,b,s)=b/s.
\end{equation}
The common scalar acts trivially, and $(\alpha,\beta)$ is jointly
surjective on rational points: $(r,b,s)=(a_0,b_0,1)$ maps to
$(a_0,b_0)$.

By \eqref{m1:eq:balanced-euler}, the Ext dimensions, with first argument
in the row and second in the column, are
\begin{equation}\label{m1:eq:decomposable-ext-table}
 \begin{array}{c|ccc}
 \dim\Ext^1&\Efac&\Rfac&\Qfac\\\hline
 \Efac&0&0&1\\
 \Rfac&0&2&1\\
 \Qfac&1&1&1
 \end{array}
\end{equation}
For example the $\Efac$--$\Rfac$ Euler value is zero, the
$\Efac$--$\Qfac$ and $\Rfac$--$\Qfac$ values are $-1$, and the
self-values are $1,-1,0$, respectively. The Hom dimensions are one
on the diagonal and zero elsewhere. Target scalar divided by source
scalar is the action on each cross-Ext. Therefore
\begin{equation}\label{m1:eq:decomposable-normal}
 N_E=F^3_{\triv}\oplus F_\alpha\oplus F_{\alpha^{-1}}
                  \oplus F_\beta\oplus F_{\beta^{-1}}.
\end{equation}
The normal space has dimension $7=16-(12-3)$.

The representative is fixed by transpose, whose ambient differential is
$(\dot A,\dot B,\dot C,\dot D)\mapsto
(\dot A^t,\dot C^t,\dot B^t,\dot D^t)$. To identify its action on
self-Ext, write the trivial ambient coordinates as
$a_1=\dot A_{11}$, $a_2=\dot A_{22}$, $b_1'=\dot B_{11}$,
$c_1'=\dot C_{11}$, $d_1=\dot D_{11}$, and $d_2=\dot D_{22}$.
For diagonal infinitesimal changes of basis,
$a_1=x(b_1'+c_1')$, $a_2$ is an independent orbit direction, and
the two $D$-coordinates are zero. The three quotient coordinates are
therefore
\begin{equation}\label{m1:eq:decomposable-fixed-coordinates}
 a_1-x(b_1'+c_1'),\qquad d_1,\qquad d_2.
\end{equation}
They are independent on the quotient by
\eqref{m1:eq:decomposable-ext-table}. Transpose interchanges $b_1',c_1'$
and fixes $a_1,d_1,d_2$, so it fixes all three.

On the stabilizer, transpose inverts $(r,b,s)$, hence both effective
characters. It exchanges each pair of cross-Ext lines. Its square is
one; choosing a rational nonzero vector in each positive line and its
image in the inverse line gives
\begin{equation}\label{m1:eq:decomposable-transpose}
 \tau_E(z_1,z_2,z_3,x,y,u,v)=(z_1,z_2,z_3,y,x,v,u).
\end{equation}
Thus the normal representation is a sum of two independently scaled
opposite-weight planes and three fixed directions, for $e\neq0$ and $xd\neq1$.

\subsubsection{The stable-system case}\label{m1:subsec:stable-normal}
At \eqref{m1:eq:stable-rep-general}--\eqref{m1:eq:H}, the decomposition is
$E=\Efac\oplus\Hfac$. Absolute stability and the distinct dimension
vectors give
\begin{equation}\label{m1:eq:stable-stabilizer}
 S=\{s(r,h)=(\diag(r,h),rI,\diag(r,h)):r,h\in F^\times\},
 \qquad\eta(r,h)=h/r.
\end{equation}
The effective character is surjective, since $\eta(1,t)=t$. This also
follows directly from cyclicity: a matrix commuting with $D$ and
satisfying $qe_1=re_1$ has $qDc=rDc$, hence is $rI$.

The balanced Euler formula now gives
\begin{equation}\label{m1:eq:stable-ext-table}
 \begin{array}{c|cc}
 \dim\Ext^1&\Efac&\Hfac\\\hline
 \Efac&0&1\\\Hfac&1&4
 \end{array},\qquad
 N_E=F^4_{\triv}\oplus F_\eta\oplus F_{\eta^{-1}}.
\end{equation}
Indeed the self-Euler values are $1,-3$ and the cross-Euler value is
$-1$. The four trivial directions are $\Ext^1(\Hfac,\Hfac)$, while the
two moving lines are cross-Ext. The total dimension is $6=16-(12-2)$.

The symmetric matrix $H$ from \eqref{m1:eq:H} satisfies
\begin{equation}\label{m1:eq:H-identities}
 HD=D^tH,\qquad He_1=b^t,\qquad E_{11}H=B,
 \qquad H^{-1}B^t=E_{11}.
\end{equation}
For instance both $HD$ and $D^tH$ equal the symmetric matrix with
entries $b_1$, $sb_1-db_0$, and $-db_1+s(sb_1-db_0)$. Consequently
\begin{equation}\label{m1:eq:stable-transporter}
 g(E)\tau(E)=E,\qquad g(E)=(I,H^{-1},I).
\end{equation}
The induced ambient operator is
\begin{equation}\label{m1:eq:stable-ambient-transpose}
 T_H(\dot A,\dot B,\dot C,\dot D)
 =(\dot A^t,\dot C^tH,H^{-1}\dot B^t,H^{-1}\dot D^tH).
\end{equation}
It preserves the orbit tangent because it fixes the base point and
normalizes $L$. Symmetry of $H$ gives $T_H^2=1$. Moreover
$\tau_Es(r,h)\tau_E^{-1}=s(r^{-1},h^{-1})$: the middle component of
$s$ is scalar and commutes with $H^{-1}$. Thus $\eta$ is inverted.
Choosing $e_{\eta^{-1}}=\tau_Ee_\eta$ makes the two cross-Ext lines an
ordinary rational coordinate swap.

It remains to identify the operator on all four self-Ext directions.
For the normalized family $E=E(a,b_0,d,s)$ its parameter derivatives are
\begin{equation}\label{m1:eq:stable-parameter-derivatives}
 \begin{aligned}
 v_a&=(E_{11},-dE_{12},0,0),&
 v_{b_0}&=(0,E_{11}+sE_{12},0,0),\\
 v_d&=(0,-aE_{12},0,-E_{12}),&
 v_s&=(0,b_0E_{12},0,E_{22}).
 \end{aligned}
\end{equation}
All have trivial weight.

\begin{lemma}\label{m1:lem:stable-parameter-basis}
Their classes form a basis of the trivial summand of the normal space,
and the corrected transpose fixes this summand pointwise.
\end{lemma}
\begin{proof}
The split stabilizer torus decomposes the complex into weight spaces,
so a weight-zero orbit tangent has a weight-zero preimage. On the
$\Hfac$-factor such a preimage consists of scalars $p,u$ at the ends
and $Q_0=(q_{ij})$ in the middle. The $\Efac$-factor has zero self-Ext.
Suppose a parameter variation is an orbit tangent. Since $c=e_1$ is
fixed, $Q_0e_1=ue_1$, so $q_{11}=u$ and $q_{21}=0$. Comparing
\[
 [Q_0,D]=\begin{pmatrix}0&-\dot d\\0&\dot s\end{pmatrix}
\]
in entries $(1,1)$ and $(2,1)$ gives $q_{12}=0$, $q_{22}=u$. Thus
$Q_0=uI$ and $\dot d=\dot s=0$. With $\delta=p-u$, the remaining
equations are
\[
 \dot a=a\delta,\qquad\dot b_0=b_0\delta,\qquad
 -d\dot a+s\dot b_0=b_1\delta.
\]
The normalization $b_1=1-ad+sb_0$ forces
\begin{equation}\label{m1:eq:stable-independence-identity}
 (-ad+sb_0-b_1)\delta=-\delta=0.
\end{equation}
All four variations vanish. The classes are independent and, by
\eqref{m1:eq:stable-ext-table}, span the four-dimensional trivial quotient.

Differentiate \eqref{m1:eq:stable-transporter} along any of these parameter
directions $v$. It gives
\[
 (dg)_E(v)\cdot\tau(E)+g(E)\tau(v)=v.
\]
The first term is an orbit tangent: with $\mathfrak l=\Lie(L)$,
$\xi=(dg)_E(v)g(E)^{-1}\in\mathfrak l$ and the term equals $\xi\cdot E$.
Hence $[g(E)\tau(v)]=[v]$ in the normal space. These four classes span
its trivial summand, proving pointwise fixation.
\end{proof}
Thus \eqref{m1:eq:stable-ext-table} has four fixed coordinates and a swap
on $F_\eta\oplus F_{\eta^{-1}}$. The argument is valid at every
$\det H\neq0$, including nilpotent and repeated-root loop values.

\subsubsection{Normal representations and transpose}\label{m1:subsec:normal-interface}
\begin{corollary}[Rational transpose stability]\label{m1:cor:transpose-stable}
Every closed $L(F)$-orbit in $\cX(F)$ is stable under $\tau$. At every
point its extended stabilizer and normal representation are as in
\eqref{m1:eq:corrected-extension} and the calculations above.
\end{corollary}
\begin{proof}
The transporters for dense and rank-zero representatives were given in
Lemma~\ref{m1:lem:dense-transpose} and \eqref{m1:eq:rankzero-transpose}.
Decomposable representatives are fixed by transpose; stable-system
representatives use \eqref{m1:eq:stable-transporter}. For $y=lE$ a
transporter is $l g(E)\alpha(l)^{-1}$, so the assertion holds on each
rational orbit. The classification is exhaustive.
\end{proof}
For the untwisted stabilizer action, two choices of transporter differ
by $L_E(F)$ and hence agree on $L_E(F)$-invariant normal distributions.

The following table summarizes the nine rational normal
representations. The last column indicates the method used in
Section~\ref{m1:sec:descendants} to prove invariance under the
corrected transpose. We write
$P_\gamma=F_\gamma\oplus F_{\gamma^{-1}}$, with its coordinate swap,
and $\mathcal V_T=\Hom(F^2,F^2)\oplus\Hom(F^2,F^2)$ with action
$(g,t)(X,Y)=(gXt^{-1},tYg^{-1})$ and the adjoint exchange
\eqref{m1:eq:rankzero-transpose}. Each $F^j$-summand in the classification is fixed by both
the stabilizer and corrected transpose. The nonsplit $K$ in the dense
row has action and odd operator \eqref{m1:eq:dense-nonsplit-action}.
\begin{center}
\small
\begin{tabular}{@{}llll@{}}
\toprule
Family & Stabilizer & Normal space & Vanishing method\\
\midrule
Dense absolute & $F^\times$ & $F^5$ & fixed operator\\
Dense split & $(F^\times)^2$ & $F^4\oplus P_\alpha$ & opposite weights\\
Dense nonsplit & $K^\times$ & $F^4\oplus K$ & norm fibers\\
Dense scalar & $\GL_2(F)$ & $M_2(F)^2$ & Prasad, translations\\
Zero split regular & $\GL_2(F)\times(F^\times)^2$ & $\mathcal V_T\oplus F^2$ & Two-Leg Lemma\\
Zero nonsplit regular & $\GL_2(F)\times K^\times$ & $\mathcal V_T\oplus F^2$ & Two-Leg Lemma\\
Zero scalar & $\GL_2(F)^2$ & $M_2(F)^2\oplus\mathfrak{gl}_2(F)$ & Identity Descendant\\
One decomposable & $(F^\times)^3$ & $F^3\oplus P_\alpha\oplus P_\beta$ & independent planes\\
One stable system & $(F^\times)^2$ & $F^4\oplus P_\eta$ & opposite weights\\
\bottomrule
\end{tabular}
\end{center}
The scalar matrix-pair row has componentwise transpose, and the scalar
rank-zero row has \eqref{m1:eq:rankzero-scalar}, including its adjoint
summand. The rational character maps are
\eqref{m1:eq:dense-split-normal}, \eqref{m1:eq:dense-nonsplit-action},
\eqref{m1:eq:decomposable-characters}, and \eqref{m1:eq:stable-stabilizer};
the remaining effective actions are given by the matrix formulas above.
In every row $\dim N_E=4+\dim\Aut(E)$: respectively
$5,6,6,8,10,10,12,7,6$. The finite extension acts on the same normal space through the corresponding
odd operator and sign character.
